% mage-auto-heading: the-capability
% mage-auto-heading: when-this-stack-earns-its-keep
% mage-auto-heading: the-composition
% mage-auto-heading: constituent-moves
% mage-auto-heading: why-these-travel-together
% Authoritative LaTeX manuscript; Markdown is a generated compatibility projection.

\section{The Capability}\label{heading:the-capability}

\textbf{Prevent actor timing from becoming accidental concurrency policy for a shared resource.} When actors contend for a scarce resource, move admission policy out of their timing and into an explicit control.

\section{When This Stack Earns Its Keep}\label{heading:when-this-stack-earns-its-keep}

Reach for it when:

\begin{itemize}
\tightlist
\item
  \textbf{Actors contend for a scarce resource} --- a test runner, a build host, a rate-limited API --- and their arrival order decides who wins.
\item
  \textbf{Availability has become an emergent property} of when jobs happen to start, rather than a stated policy.
\item
  \textbf{Capacity should flex under live pressure}, but only once a fixed policy and a single admission point already exist.
\end{itemize}

\section{The Composition}\label{heading:the-composition}

\begin{figure}[htbp]
\centering
\includegraphics[width=\linewidth]{book/latex/assets/resource-mediation-stack-7443deee0a.pdf}
\caption{The resource-mediation composition. A RESOURCE POLICY states which resource is scarce and what capacity is acceptable; a MEDIATOR admits at most N actors through one admission point to the SHARED RESOURCE. A dashed LIVE PRESSURE loop tightens or relaxes effective capacity when justified. N=1 and N\textgreater1 are settings of the same mediator, not separate stacks. Solid path: the required path. Dashed attachment: a useful enhancement, not required for the capability.}
\label{resource-mediation-stack}
\end{figure}

\section{Constituent Moves}\label{heading:constituent-moves}

{\footnotesize
{\def\LTcaptype{none} % do not increment counter
\begin{longtable}[]{@{}
  >{\raggedright\arraybackslash}p{(\linewidth - 2\tabcolsep) * \real{0.5000}}
  >{\raggedright\arraybackslash}p{(\linewidth - 2\tabcolsep) * \real{0.5000}}@{}}
\toprule\noalign{}
\begin{minipage}[b]{\linewidth}\raggedright
Move
\end{minipage} & \begin{minipage}[b]{\linewidth}\raggedright
Role
\end{minipage} \\
\midrule\noalign{}
\endhead
\bottomrule\noalign{}
\endlastfoot
\textbf{DECLARE} & State which resource is scarce and what capacity is acceptable. \\
\textbf{MEDIATE} & Route access through one admission point. \\
\textbf{ADAPT} & \emph{(strengthens)} Modify effective capacity from live pressure where justified. \\
\end{longtable}
}
}

\section{Why These Travel Together}\label{heading:why-these-travel-together}

Unmediated contention makes resource availability an emergent property of actor timing---whoever starts first wins, and the policy is an accident. A mediator replaces that accident with explicit policy: one admission point and a declared capacity. The engineering move is not the semaphore; it is moving contention policy out of individual actors and into one shared control. Whether that control admits one actor or many is a setting, not a different mechanism.

Adaptive control can improve utilization, but only after a fixed policy and common admission point exist. Capacity cannot be tuned coherently when no single control owns it.

\textbf{Mechanisms:} exclusive mediator · bounded mediator · pressure-responsive admission
